\documentclass[11pt]{amsart}

\usepackage[a4paper,margin=1.02in]{geometry}
\usepackage{amsmath,amssymb,amsthm,mathtools,mathrsfs,bm}
\usepackage{booktabs,tabularx,array,longtable}
\usepackage{enumitem}
\usepackage{tikz-cd}
\usepackage{microtype}
\usepackage{xcolor}
\usepackage{hyperref}
\usepackage[nameinlink,capitalize,noabbrev]{cleveref}
\usepackage[most]{tcolorbox}
\usepackage{listings}

\hypersetup{
  colorlinks=true,
  linkcolor=blue!45!black,
  citecolor=blue!45!black,
  urlcolor=blue!45!black,
  pdftitle={A Comprehensive Reconstruction of the IUT-ABC Proof Programme},
  pdfauthor={ChatGPT}
}

\newtheorem{theorem}{Theorem}[section]
\newtheorem{proposition}[theorem]{Proposition}
\newtheorem{lemma}[theorem]{Lemma}
\newtheorem{corollary}[theorem]{Corollary}
\newtheorem{counterexample}[theorem]{Counterexample}
\theoremstyle{definition}
\newtheorem{definition}[theorem]{Definition}
\newtheorem{construction}[theorem]{Construction}
\theoremstyle{remark}
\newtheorem{remark}[theorem]{Remark}
\newtheorem{status}[theorem]{Status}

\newcommand{\R}{\mathbb R}
\newcommand{\Q}{\mathbb Q}
\newcommand{\PP}{\mathbb P}
\newcommand{\OO}{\mathcal O}
\newcommand{\UTheta}{\mathcal U_\Theta}
\newcommand{\mup}{\mu_{\mathrm{proc}}}
\newcommand{\rad}{\operatorname{rad}}
\newcommand{\Reg}{\operatorname{Reg}}
\newcommand{\eps}{\varepsilon}
\newcommand{\FMod}{\mathcal F_{\mathrm{MOD}}}
\newcommand{\Fmod}{\mathcal F_{\mathrm{mod}}}

\title[Comprehensive IUT--$abc$ reconstruction]{A Comprehensive Reconstruction of the IUT--$abc$ Proof Programme\\
\large Proven components, countermodels, normalization corrections, non-circular closure, and the exact remaining construction theorems}
\author{ChatGPT}
\date{August 20, 2026}

\subjclass[2020]{Primary 11D75, 11G50; Secondary 03B35, 18A40, 28C10}
\keywords{ABC conjecture, inter-universal Teichm\"uller theory, Lean, poly-isomorphism, q-pilot, log-volume, formal verification}

\begin{document}
\maketitle

\begin{abstract}
This paper consolidates the mathematical and formal results obtained in a
long-form reconstruction of the proposed inter-universal Teichm\"uller
(IUT) route to the $abc$ conjecture.  The work separates three layers that
must not be conflated: elementary and categorical implications that can be
proved independently of the full IUT source; source-level reconstructions of
the pilot bridge; and the actual inhabitation theorems required to turn the
public Lean specification into a parameter-free proof.

The following components are proved in ordinary mathematics and encoded in
Lean source: representative extraction from equality of composite
poly-morphisms; Kummer conjugation for full poly-isomorphisms; a finite
countermodel showing that projection-wise compatibility is insufficient;
generation of a theta-pilot region as the union of concrete outputs;
finite-positive log-volume monotonicity; ramification-correct q-pilot
normalization; product-weight marginalization; procession averaging;
finite-exception absorption; the scalar IUT IV algebra; the geometric tripod
to logarithmic $abc$ transfer; the correct pointwise quantifier structure;
and a circularity theorem showing that a former downstream certificate
already contained the $abc$ conjecture as a field.

A non-circular final certificate is defined.  Its inhabitation would produce
a parameter-free Lean theorem stating the standard logarithmic $abc$
conjecture.  No inhabitant is constructed here.  The remaining obligations
are an actual initial-theta-data family, an actual IUT III
Hodge-theater/Frobenioid/multiradial source family, and a uniform IUT IV
geometric height bridge.  The public IUT formalization deliberately leaves
these objects behind interfaces.  Consequently this paper is a comprehensive
reconstruction and formalization audit, not an unconditional proof of $abc$.
\end{abstract}

\begin{tcolorbox}[colback=yellow!5,colframe=black!65,title=Final status]
The paper proves a large collection of genuine lemmas and an honest
non-circular closure theorem
\[
  \texttt{FourStageProgram}\Longrightarrow\texttt{ABCConjecture}.
\]
It does not prove
\[
  \operatorname{Nonempty}(\texttt{FourStageProgram}).
\]
Thus no theorem of the form
\[
  \texttt{theorem abc\_conjecture : ABCConjecture}
\]
has been obtained without parameters.
\end{tcolorbox}

\tableofcontents

\section{The target statement}

For a positive integer $n$, let
\[
 \rad(n)=\prod_{p\mid n}p
\]
be the product of the distinct prime divisors.  The logarithmic form of the
$abc$ conjecture is the following.

\begin{theorem}[ABC conjecture --- target statement]\label{thm:abc-target}
For every $\eps>0$ there exists $C_\eps\in\R$ such that, for all positive
pairwise coprime integers $a,b,c$ with $a+b=c$,
\[
 \log\max\{a,b,c\}
 \le
 (1+\eps)\log\rad(abc)+C_\eps.
\]
\end{theorem}

The formal development defines exactly this proposition.  The issue is not the
syntax of the final statement.  The issue is constructing, uniformly in the
arithmetic input and in $\eps$, all source and height-theoretic data needed to
prove it.

\section{The IUT numerical corridor}

The decisive IUT III corridor is usually displayed as
\[
\begin{tikzcd}[column sep=small]
W\in\UTheta,\quad \mup(W)=-|\log q|
\arrow[r]
&
-|\log q|\le-|\log\Theta|
\arrow[r]
&
C_\Theta\ge-1.
\end{tikzcd}
\]
The first implication uses monotonicity of the procession-normalized volume;
the second combines the bridge with
\[
 -|\log\Theta|\le C_\Theta|\log q|,
 \qquad |\log q|>0.
\]
The downstream IUT IV argument is intended to turn $C_\Theta\ge-1$ into a
Vojta-type inequality on the tripod
$\PP^1\setminus\{0,1,\infty\}$ and then into \cref{thm:abc-target}.

The research programme therefore decomposes into four stages:

\begin{enumerate}[label=(\arabic*)]
\item construct an actual family of initial theta-data;
\item construct an actual IUT III multiradial output family and the witness
$W$;
\item prove a uniform, non-circular IUT IV geometric bridge;
\item compose the preceding stages to obtain the parameter-free $abc$
theorem.
\end{enumerate}

\section{Weak compatibility is not enough}

The first important negative result is that a common source, a horizontal
pilot map, compatibility on projected unit data, and the presence of an
ordinary Kummer branch do not by themselves force a coherent pilot output.

\begin{counterexample}[Finite common-source countermodel]\label{cex:weak}
Let
\[
 A_0=A_1=C_0=C_1=\{0,1\}.
\]
Take the theta-pilot to be $0$, the q-pilot to be $1$, and let the horizontal
map interchange $0$ and $1$.  Let both Kummer maps be identities.  Suppose
the only allowed coric transport is the identity.  Make every unit and label
projection a singleton, let Ind2 act trivially, and include an ordinary branch
whose possible output is $0$.

Then all projected compatibility conditions hold and the horizontal map sends
the theta-pilot to the q-pilot, but no allowed coric map sends the Kummer image
of $0$ to the Kummer image of $1$.
\end{counterexample}

This counterexample does not refute IUT.  It refutes the inference from a list
of weak projection-wise properties to the existence of one object-level
coherent tuple.  It also explains why a formalization must identify the exact
level at which fullness or compatibility is asserted.

\section{Poly-morphism representative extraction}

Let
\[
 K_0:A_0\overset\sim\longrightarrow C_0,
 \qquad
 K_1:A_1\overset\sim\longrightarrow C_1
\]
be Kummer isomorphisms.  Let
\[
 H\subseteq\operatorname{Iso}(A_0,A_1),
 \qquad
 P\subseteq\operatorname{Iso}(C_0,C_1)
\]
be poly-isomorphisms.  Define
\[
 K_1H=\{K_1\circ h:h\in H\},
 \qquad
 PK_0=\{p\circ K_0:p\in P\}.
\]

\begin{theorem}[Representative extraction]\label{thm:extract}
If
\[
 K_1H=PK_0,
\]
then for every $h\in H$ there exists $p\in P$ such that
\[
 pK_0=K_1h.
\]
\end{theorem}

\begin{proof}
The composite $K_1h$ lies in $K_1H$.  By equality of the two composite sets,
it lies in $PK_0$, hence equals $pK_0$ for some $p\in P$.
\end{proof}

This elementary theorem was one of the first points to be formalized in Lean.
It isolates the precise logical difference between compatibility of complete
composite poly-morphisms and separate compatibility of their projections.

\section{Full-poly conjugation}

If the coric relation is the full poly-isomorphism, the explicit equality of
composite sets is unnecessary.

\begin{definition}[Kummer conjugate]
Given a horizontal isomorphism $h:A_0\simeq A_1$, define
\[
 P=K_1hK_0^{-1}:C_0\simeq C_1.
\]
\end{definition}

\begin{theorem}[Full-poly conjugation]\label{thm:conjugate}
For every $x\in A_0$,
\[
 P(K_0x)=K_1(hx).
\]
If the allowed coric relation is the full poly-isomorphism
$\operatorname{Iso}(C_0,C_1)$, then $P$ is automatically an allowed
representative.
\end{theorem}

\begin{proof}
The first assertion is cancellation.  The second is the definition of a full
poly-isomorphism.
\end{proof}

The Boolean countermodel in \cref{cex:weak} shows that fullness is essential:
for a proper restricted family the conjugate need not be allowed.

\section{Fixed-branch reconstruction}

The actual possible-image set is a union over many choices.  To prove
existence of one witness, it is not necessary to control all branches
simultaneously.

Fix a horizontal representative $h$, a Kummer input index, the identity
representatives for Ind1 and Ind2, and the ordinary Kummer branch for Ind3.
The source statements used by this reconstruction are:

\begin{enumerate}[label=(\roman*)]
\item horizontal LGP/lgp links map theta-pilot objects to q-pilot objects;
\item the local upper-semicompatible containers contain the ordinary Kummer
image;
\item local fractional ideals globalize to objects of the corresponding
global Frobenioids;
\item global log-volume agrees with arithmetic degree after normalization.
\end{enumerate}

Under the structured compatibility hypothesis, the fixed branch produces a
global fractional-ideal region $W$ satisfying
\[
 W\in\UTheta,
 \qquad
 \mup(W)=-|\log q|.
\]
The numerical bridge follows.

This is a source-level reconstruction.  It is not yet an object generated by a
complete Lean encoding of Hodge theaters and Frobenioids.

\section{Generated possible-image unions}

An independent membership field can be removed by defining the theta-pilot
region from the concrete outputs.

Let $\mathscr O$ be a nonempty type of outputs and let $R_{o,i}$ be the
admissible region associated to output $o$ in capsule $i$.  Assume that for
each $i$ there is a finite set $S_i$ of rational places such that
\[
 v_{\Q}\notin S_i
 \Longrightarrow
 R_{o,i}(v_{\Q})=\OO_{i,v_{\Q}}
\]
for all outputs $o$.

\begin{definition}[Generated theta-pilot union]
Define
\[
 T_i(v_{\Q})=\bigcup_{o\in\mathscr O}R_{o,i}(v_{\Q}).
\]
\end{definition}

\begin{theorem}[Admissibility and definitional membership]\label{thm:union}
Each $T_i$ is a restricted-product admissible region, and
\[
 R_{o,i}\subseteq T_i
\]
for every $o$.
\end{theorem}

\begin{proof}
Outside $S_i$ all regions in the union equal the same integral region.  Since
$\mathscr O$ is nonempty, the union equals that integral region there.  Hence
the support is finite.  Membership is the canonical injection into a union.
\end{proof}

This construction was implemented directly against the public
\texttt{LargeVolumeContainerData.AdmissibleRegion} type.  It removes the
separate field asserting that the native output belongs to a previously
supplied theta region.

\section{Actual public-region witness theorem}

Let
\[
 X:\texttt{Iut.Corollary312VariantData}.
\]
An actual witness consists of a capsule-indexed family of public admissible
regions $W_i$, containment $W_i\subseteq\Theta_i$, an exact q-volume identity,
and monotonicity on the two good region families used in the argument.

\begin{theorem}[Actual witness closes the public numerical variant]
An actual public-type witness implies
\[
 \texttt{Iut.Corollary312Variant}(X).
\]
\end{theorem}

\begin{proof}
The public hull system gives
\[
 \Theta_i\subseteq H(\Theta_i).
\]
Thus $W_i\subseteq H(\Theta_i)$.  Using the exact q-volume identity and
procession-volume monotonicity gives
\[
 X.\texttt{qPilot.lhs}
 \le
 X.\texttt{rhsData.rhs}.
\]
\end{proof}

This is one of the strongest formal results obtained: the numerical theorem no
longer uses an auxiliary region type.  The remaining problem is constructing
the actual witness from complete IUT source definitions.

\section{Finite-positive log-volume monotonicity}

The public volume interface is total on arbitrary regions and therefore does
not impose a global monotonicity axiom.  For finite positive measure-realized
regions, the required theorem is standard.

\begin{theorem}[Logarithmic measure monotonicity]\label{thm:log-mono}
Let $\mu$ be a measure and let $A\subseteq B$.  If
\[
 \mu(A)>0,
 \qquad
 \mu(B)<\infty,
\]
then
\[
 \log\mu(A)\le\log\mu(B).
\]
\end{theorem}

\begin{proof}
Measure monotonicity gives $\mu(A)\le\mu(B)$.  The real logarithm is
increasing on $(0,\infty)$.
\end{proof}

Nonnegative finite weighted sums and finite averages preserve pointwise
inequalities.  Hence packet and procession monotonicity follow once the actual
component log-volumes are represented by finite positive Haar measures.

\section{Native q-pilot normalization}

Let
\[
 D_q=\sum_v n_v[v]
\]
be the arithmetic q-divisor, where $n_v=\operatorname{ord}_v(q_v)$ is the
integer uniformizer order.  Define
\[
 \log(q)=\deg_{\mathrm{norm}}(D_q).
\]
The q-pilot uses a $2\ell$-th root of the Tate parameter, so its realified
divisor is
\[
 D_{\mathbf q}=\frac1{2\ell}D_q.
\]
By linearity,
\[
 \deg_{\mathrm{norm}}(D_{\mathbf q})
 =\frac1{2\ell}\log(q)=|\log q|.
\]
With the signed volume convention,
\[
 \boxed{\mup(Q_{\mathrm{pilot}})=-|\log q|.}
\]
This closes the scalar calibration of the native object.  It does not prove
that the same object is an actual member of the theta-pilot possible-image
union.

\section{Ramification correction}

Let $L/\Q_p$ have ramification index $e$ and residue degree $f$.  Write
\[
 p=u\pi^e,
 \qquad
 q=u_q\pi^n,
\]
with units $u,u_q$.  Suppose $\rho$ is the packet-normalized logarithmic
response, additive under multiplication, zero on units, and normalized by
\[
 \rho(p)=-\log p.
\]
Then
\[
 \rho(\pi)=-\frac1e\log p,
 \qquad
 \rho(q)=-\frac ne\log p.
\]
Multiplying by the normalized local-degree weight
\[
 \frac{ef}{[K:\Q]}
\]
cancels $e$ and yields
\[
 -\frac{nf}{[K:\Q]}\log p,
\]
which is the normalized arithmetic-divisor contribution.

This proves a normalization fork:

\begin{itemize}
\item local-degree weight $ef/d$ must be paired with p-normalized order $n/e$;
\item integer uniformizer order $n$ must be paired with residue-degree weight
$f/d$.
\end{itemize}

Using local-degree weight and integer order simultaneously introduces an
extra factor $e$.  A ramified quadratic counterexample was formalized.

\section{Packet marginalization and procession averaging}

Let $S$ be a finite label set, let $V_p$ be the finite fiber over $p$, and let
$w:V_p\to\R$ satisfy
\[
 \sum_{v\in V_p}w(v)=1.
\]
A packet component is a function $c:S\to V_p$ with product weight
\[
 W(c)=\prod_{j\in S}w(c(j)).
\]

\begin{theorem}[Product-weight marginalization]\label{thm:marginal}
For every distinguished label $j_0\in S$ and every $f:V_p\to\R$,
\[
 \sum_{c:S\to V_p}W(c)f(c(j_0))
 =
 \sum_{v\in V_p}w(v)f(v).
\]
\end{theorem}

\begin{proof}
Split $c$ into its $j_0$ coordinate and the remaining coordinates.  For a
fixed value of the distinguished coordinate, the sum of the weights over the
remaining coordinates is the product of copies of $\sum_vw(v)=1$.
\end{proof}

More generally, any barycentric allocation across labels whose coefficients
sum to one has the same normalized reading.  A naive diagonal action at every
label is impossible: it multiplies a nonzero response by the capsule
cardinality, which varies along the procession.

If every capsule has the same calibrated global q-reading, its procession
average is that same value.  This removes procession averaging as an
independent hypothesis.

\section{Holomorphic hull reductions}

Two independent reductions were obtained.

First, if a generated union lies inside a log-shell whose closure is compact,
then the closure of the generated union is compact.  This proves the relative
compactness portion of hull admissibility.

Second, suppose a family of regions is indexed by exponents in an antitone
chain.  If the exponent family has an attained lower bound $m$, then the
region at exponent $m$ is the least member of the chain containing the union.
This proves the order-theoretic portion of a least-hull construction.

The remaining local-field task is to instantiate the chain by the actual
nonarchimedean lattices and the archimedean closed balls or polydiscs in the
IUT packet summands.

\section{The IUT IV downstream reconstruction}

Once the corrected Corollary 3.12 estimate is available, the final scalar
algebra is elementary.  A typical coefficient has the form
\[
 C_\Theta
 =\lambda\left(A-\frac16\log q\right)-1,
 \qquad \lambda>0.
\]
Then $C_\Theta\ge-1$ implies
\[
 \frac16\log q\le A.
\]
The subsequent $\eps$-absorption inequalities have also been formalized.

The geometric part has three standard components.

\subsection{The Legendre height comparison}
For
\[
 E_\lambda:y^2=x(x-1)(x-\lambda),
\]
the $j$-invariant is
\[
 j(\lambda)=2^8\frac{(1-\lambda+\lambda^2)^3}
 {\lambda^2(1-\lambda)^2}.
\]
Its pole divisor is $2[0]+2[1]+2[\infty]$.  Hence
\[
 h_\infty(j(\lambda))
 =6h_{\omega_{\PP^1}(0+1+\infty)}(\lambda)+O(1).
\]
Tate uniformization identifies $\log q$ with the pole-height of $j$ up to
bounded discrepancy on the relevant compactly bounded family, yielding
\[
 \frac16\log q
 =h_{\omega_{\PP^1}(0+1+\infty)}+O(1).
\]

\subsection{Tripod different and conductor}
For a coprime triple $a+b=c$ and $\lambda=a/c$,
\[
 h(\lambda)=\log\max\{a,b,c\},
\]
the relative different over $\Q$ is zero, and
\[
 \log\operatorname{Cond}_{\{0,1,\infty\}}(\lambda)
 =\sum_{p\mid abc}\log p
 =\log\rad(abc).
\]
Thus the tripod Vojta inequality specializes exactly to the logarithmic
$abc$ inequality.

\subsection{Exceptional sets and admissible primes}
A finite exceptional set can be absorbed into the bounded-discrepancy
constant.  The elementary prime-selection step --- choosing a prime above a
bound and outside a finite set --- is also formalized.  The non-elementary
part is proving that the Galois-image, Tate-order, and initial-theta
conditions fail only on an admissible exceptional family; that remains part
of the actual IUT IV source theorem.

\section{Tripod Statement II implies abc}

Let $T$ be an actual geometric tripod height theory.  Suppose there is an
encoding of $abc$ triples into points of $T$ and constants $A_h,A_d,A_c$ such
that
\[
 \log\max\{a,b,c\}
 \le h_T(P)+A_h,
\]
\[
 \log\operatorname{Diff}_T(P)\le A_d,
\]
\[
 \log\operatorname{Cond}_T(P)
 \le\log\rad(abc)+A_c.
\]
If Statement II gives, uniformly in $P$,
\[
 h_T(P)
 \le
 (1+\eps)
 \bigl(\log\operatorname{Diff}_T(P)+
       \log\operatorname{Cond}_T(P)\bigr)+B_\eps,
\]
then the logarithmic $abc$ conjecture follows after absorbing the comparison
constants.  This implication is formalized without storing $abc$ as a field.

\section{The quantifier correction}

A universal Diophantine theorem cannot follow from one isolated initial
Theta datum.  The correct source is a dependent family
\[
 x\longmapsto
 \bigl(D_x,Q_x,S_x\bigr)
\]
indexed by all arithmetic inputs to which the theorem applies.

A finite Boolean countermodel formalizes the logical point:
\[
 \exists x\,\operatorname{Certificate}(x)
 \not\Rightarrow
 \forall x\,\operatorname{Certificate}(x).
\]

Thus the first two construction goals must have the form
\[
 \forall x:\mathrm{Input},\quad
 D_x:\texttt{InitialThetaData},
\]
\[
 \forall x:\mathrm{Input},\quad
 S_x:\texttt{GeneratedNativeSource}(D_x,Q_x).
\]

\section{Circularity of the former downstream certificate}

An earlier certificate used a field
\[
 \bigl(\forall x,\operatorname{Cor312}(x)\bigr)
 \longrightarrow ABCConjecture.
\]
For a fixed source family $F$, the type of such downstream objects is
inhabited if and only if the target proposition is already true:
\[
 \boxed{
 \operatorname{Nonempty}
 (\operatorname{UniformDownstream}(F,T))
 \iff T.
 }
\]
Consequently the former three-closure certificate was inhabited precisely
when an upstream family existed and $abc$ was already true.  This circularity
was proved in Lean.

\section{The non-circular final certificate}

The circular field is replaced by explicit intermediate data:

\begin{itemize}
\item an encoding of $abc$ points into the input family;
\item functions $\log Q$, $\log\mathrm{Diff}$, and $\log\mathrm{Cond}$;
\item uniform height, different, and conductor comparison constants;
\item for each $\eps>0$, one constant $C_\eps$ such that Corollary 3.12 for
the encoded input implies the uniform q-height estimate.
\end{itemize}

This structure is called a non-circular IUT IV bridge.  Lean proves that a
pointwise actual IUT III source family together with this bridge implies the
standard logarithmic $abc$ conjecture.

\section{The four-stage programme}

The final honest formal structure has three construction stages and one
consequence:

\begin{enumerate}[label=(\arabic*)]
\item an actual initial Theta-data family;
\item an actual IUT III multiradial source family;
\item an actual non-circular uniform IUT IV geometric bridge;
\item the theorem $abc$.
\end{enumerate}

The Lean project defines a structure \texttt{FourStageProgram} containing the
first three stages and proves
\[
 \boxed{
 \texttt{FourStageProgram}
 \Longrightarrow
 \texttt{ABCConjecture}.
 }
\]
Equivalently,
\[
 \operatorname{Nonempty}(\texttt{FourStageProgram})
 \Longrightarrow
 ABCConjecture.
\]

No inhabitant of \texttt{FourStageProgram} is constructed.

\section{Why the four missing constructions were not solved here}

The remaining problem is not a shortage of wrapper definitions.  It is the
absence of proofs inhabiting the fields.

\subsection{Actual initial Theta-data family}
The public structure contains a number field, an elliptic curve, an admissible
prime, a genuine mod-$\ell$ representation, a torsion field, orbicurves,
coverings, fundamental groups, core relations, cusps, Tate parameters, and
local conditions.  Defining a record with these field names does not construct
these objects or prove their compatibility.

\subsection{Actual IUT III source family}
The public Corollary 3.12 project takes the theta region, packet presentation,
log-volume, and hull system as interface data.  It does not define a complete
Hodge-theater/Frobenioid/log-link/Kummer/multiradial source whose output
inhabits those interfaces.  The generated-union construction proves what to
do with concrete outputs; it does not manufacture the outputs.

\subsection{Actual IUT IV uniform bridge}
The non-circular bridge requires one uniform constant for all arithmetic
points and all the actual height, different, and conductor comparisons.  It
also requires the source theorems controlling exceptional sets, admissible
primes, and the actual local/global packet estimates.  Packaging these
statements as fields does not prove them.

\subsection{Parameter-free abc theorem}
The final theorem is a consequence of the preceding stages, not an additional
independent construction.  A declaration
\[
 \texttt{theorem abc\_conjecture : ABCConjecture}
\]
without a term inhabiting the first three stages would hide an unproved axiom.

\section{Ledger of results}

\begin{longtable}{@{}p{0.46\textwidth}p{0.19\textwidth}p{0.28\textwidth}@{}}
\toprule
Result & Status & Comment\\
\midrule
\endfirsthead
\toprule
Result & Status & Comment\\
\midrule
\endhead
Standard logarithmic $abc$ statement & Formalized & Exact target proposition\\
Weak common-source sufficiency & Refuted & Finite Boolean countermodel\\
Composite poly-morphism extraction & Proved & Elementary set-valued composition\\
Full-poly Kummer conjugation & Proved & Requires fullness at the correct level\\
Ordinary generated membership & Proved & Definitional in generated choice type\\
Actual public witness $\Rightarrow$ public Cor312 variant & Proved & Uses public region/hull types\\
Generated theta union admissible & Proved conditionally & Requires common finite support\\
Finite-positive log-volume monotonicity & Proved & Standard measure theory\\
Native q-pilot degree calibration & Proved abstractly & Uses linear arithmetic degree\\
Ramification correction & Proved & Local response is $-(n/e)\log p$\\
Packet marginalization & Proved & Finite product-measure identity\\
Procession averaging & Proved & Constant capsule reading\\
Finite exceptional-set absorption & Proved & Does not prove Galois-finiteness\\
Elementary prime avoidance & Proved & Does not prove large Galois image\\
IUT IV scalar algebra & Proved & Does not prove actual packet estimates\\
Tripod Statement II $\Rightarrow abc$ & Proved conditionally & Requires actual comparison package\\
Pointwise quantifier correction & Proved & One certificate is insufficient\\
Circularity of old downstream certificate & Proved & Old certificate already contained $abc$\\
Non-circular bridge $\Rightarrow abc$ & Proved & Honest final logical closure\\
Actual initial Theta-data family & Open & No source inhabitant constructed\\
Actual IUT III source family & Open & Complete source types absent\\
Actual uniform IUT IV bridge & Open & Uniform geometric theorem absent\\
Parameter-free $abc$ theorem & Open & Equivalent to inhabiting prior stages\\
\bottomrule
\end{longtable}

\section{Lean package architecture}

The accompanying project contains thirty-seven Lean modules.  Its main groups
are:

\begin{description}[leftmargin=0pt,labelindent=0pt,style=nextline,itemsep=0.45em]
\item[Public Corollary 3.12 adapter]
\texttt{ActualPilotWitness.lean}, \texttt{GeneratedSource.lean}.
\item[Quantifier and circularity audit]
\texttt{QuantifierCorrectClosure.lean}, \texttt{CircularityAudit.lean},
\texttt{NonCircularDownstream.lean}.
\item[Full-poly and countermodel core]
\texttt{FullPolyCore.lean}, \texttt{WeakCompatibilityCountermodel.lean}.
\item[Q-pilot normalization]
\texttt{QPilotNormalizationAudit.lean},
\texttt{QPilotNormalizationFork.lean},
\texttt{RootQPilotDivisor.lean},
\texttt{RamificationCorrectedQPilot.lean}.
\item[Packets, volumes, and hulls]
\texttt{FinitePositiveLogVolumeMonotonicity.lean},
\texttt{ProductWeightMarginalization.lean},
\texttt{PrimePowerQPilotRegion.lean},
\texttt{NativeQPilotCalibration.lean},
\texttt{FiniteExponentHull.lean}.
\item[Downstream algebra]
\texttt{FiniteExceptionalSet.lean},
\texttt{AdmissiblePrimeSelection.lean},
\texttt{Cor312CoefficientAlgebra.lean},
\texttt{IUTIVAbsorption.lean}.
\item[Final honest closure]
\texttt{FourOpenConstructions.lean},
\texttt{AxiomAudit.lean}.
\end{description}

The source contains no explicit \texttt{sorry}, \texttt{admit}, or
user-declared \texttt{axiom}.  The current runtime did not contain Lean, so the
v4.0 archive received static auditing but not a fresh local kernel build.

\section{Conclusion}

The research programme has reached a sharp boundary.  The categorical,
measure-theoretic, normalization, averaging, scalar-algebra, quantifier, and
final-transfer components are no longer vague.  They have been stated as
precise theorems, many of them formalized.

What remains is substantially deeper: actual construction of the source
objects that the public Lean interfaces deliberately leave abstract.  No
amount of renaming certificates can replace the existence theorems for those
objects.  Accordingly, the present work is a comprehensive reconstruction and
formalization audit, but it does not claim an unconditional solution of the
$abc$ conjecture.

\appendix

\section{Exact non-circular theorem}

The formal theorem has the schematic shape
\[
\begin{aligned}
&\text{pointwise actual IUT III source family}\\
&\quad+\text{uniform q-height estimate from Corollary 3.12}\\
&\quad+\text{height/different/conductor comparisons}\\
&\qquad\Longrightarrow ABCConjecture.
\end{aligned}
\]
The exact remaining proposition is
\[
 \operatorname{Nonempty}(\texttt{FourStageProgram}).
\]

\section{Bibliographic sources}

\begin{thebibliography}{99}

\bibitem{IUT1}
S.~Mochizuki,
\emph{Inter-universal Teichm\"uller Theory I: Construction of Hodge Theaters},
Publ. RIMS 57 (2021), 3--207.

\bibitem{IUT2}
S.~Mochizuki,
\emph{Inter-universal Teichm\"uller Theory II: Hodge--Arakelov-theoretic Evaluation},
Publ. RIMS 57 (2021), 209--401.

\bibitem{IUT3}
S.~Mochizuki,
\emph{Inter-universal Teichm\"uller Theory III: Canonical Splittings of the
Log-theta-lattice},
Publ. RIMS 57 (2021), 403--626.

\bibitem{IUT4}
S.~Mochizuki,
\emph{Inter-universal Teichm\"uller Theory IV: Log-volume Computations and
Set-theoretic Foundations},
Publ. RIMS 57 (2021), 627--723.

\bibitem{GenEll}
S.~Mochizuki,
\emph{Arithmetic Elliptic Curves in General Position}.

\bibitem{LANA-IUT}
LANA contributors,
\emph{Inter-universal Teichm\"uller theory: the ABC/IUT trunk},
Lean 4 formalization project.

\end{thebibliography}

\end{document}
